% !TeX root = ../main.tex
\section{Short-strip detector}
\label{sec:short-strips}
The intermediate tracker uses rectangular two-coordinate strixel sensors in four
azimuthally tilted barrel layers and six petal discs at each end. The selected
working simulation model extends the short-strip endcap active coverage and
uses a common redesigned module and power distribution. Fine local $u$ is
transverse to the radial direction; coarse local $v$ follows $z$ in the barrel
and points radially inward on the discs. These dimensions are nODD design
choices. Electronics, thermal contacts and services remain engineering design
allocations, with their qualification requirements stated below.

\subsection{Sensor and redesigned module}
\Cref{tab:short-module} specifies the sensitive element and its passive stack.
The cell grid contains $640\times48=\num{30720}$ addresses per sensor. A
\qty{0.5}{\milli\metre} guard surrounds the active rectangle; two
\qty{3}{\milli\metre} flex ledges give an occupied outline of
\qty{49}{\milli\metre} by \qty{103}{\milli\metre}.

\begin{table}[!htbp]
\centering\small
\caption{Short-strip module representation. The readout tiles are logical
footprint and heat-source allocations, not an identified qualified ASIC.}
\label{tab:short-module}
\begin{tabular}{@{}lp{87mm}@{}}\toprule
Parameter & Selected representation\\\midrule
Active silicon & $48\times96\times0.20\,\unit{\milli\metre^3}$\\
Readout pitch & $75\,\unit{\micro\metre}\times2\,\unit{\milli\metre}$\\
ASIC allocation & Two $23.4\times23.4\times0.15\,\unit{\milli\metre^3}$ tiles at $(u,v)=(0,\pm18)\,\unit{\milli\metre}$\\
Interconnect & \qty{25}{\micro\metre} bump envelope and \qty{50}{\micro\metre} redistribution layer\\
ASIC bond & \qty{50}{\micro\metre} adhesive; required cold conductivity $\geq1.2\,\unit{\watt\per\metre\per\kelvin}$\\
Spreader & \qty{0.35}{\milli\metre} graphite; required $k_\parallel\geq400$, $k_\perp\geq10\,\unit{\watt\per\metre\per\kelvin}$\\
Thermal isolation & \qty{0.10}{\milli\metre} electrically insulating TIM, required cold $k\geq1.2\,\unit{\watt\per\metre\per\kelvin}$\\
Backing & \qty{0.20}{\milli\metre} CFRP, with explicit thermal windows\\
Thermal contacts & Barrel: two $20\times64\,\unit{\milli\metre^2}$ lands at $u=\pm12\,\unit{\milli\metre}$; endcap: one central $16\times60\,\unit{\milli\metre^2}$ saddle\\
Pipe interface & Direct conductive inserts and shaped pipe contacts; heat-flow path $\leq8\,\unit{\milli\metre}$, effective $k\geq100\,\unit{\watt\per\metre\per\kelvin}$\\\bottomrule
\end{tabular}
\end{table}

The direct inserts pass through windows in the backing and cold structure to
contact the titanium pipe; displaced backing, skin and core volumes are
subtracted explicitly. The pipe's outer envelope is cut from the saddle, while
coolant remains a separate volume. This represents a changed module and carrier
interface. It requires new sensor-to-ASIC redistribution, electrical isolation,
mounts and strain relief. The two tiles each allocate \num{15360} channels;
that arithmetic establishes neither an available ASIC nor a feasible bump map.
Readout power, longer-cell capacitance, noise, radiation lifetime and bonding
must be demonstrated by a physical module design.

\subsection{Tilted barrel and cold staves}
The anchor radii are $260,340,480,660\,\unit{\milli\metre}$, with
$44,56,80,108$ staves respectively. Each carries 28 sensors, giving 288 staves
and \num{8064} sensors. Endpoint sensor centres are at
$z=\pm\qty{1152}{\milli\metre}$, with \qty{85.333}{\milli\metre} row pitch
and \qty{10.667}{\milli\metre} axial active overlap. Alternating rows are
raised by \qty{1.5}{\milli\metre}; the active edges reach
$|z|=\qty{1200}{\milli\metre}$. The complete stave has a common
$+\qty{15}{\degree}$ azimuthal tilt.

\begin{figure}[tbp]
\centering\includegraphics[width=\linewidth]{short-strip-barrel}
\caption{Short-strip barrel populations and an innermost cold-stack projection.
Different longitudinal rows are superposed; the right panel is not a physical
cut at one $z$. The sensor, cold structure, tubes and local services share the
\qty{15}{\degree} tilt.}
\label{fig:short-barrel}
\end{figure}

The \qty{52}{\milli\metre}-wide cold structure reaches
$|z|=\qty{1210}{\milli\metre}$. Its sandwich contains a
\qty{5}{\milli\metre} conductive carbon-foam core,
\qty{0.10}{\milli\metre} adhesive and \qty{0.15}{\milli\metre} CFRP skins on
each face. Nine CFRP bearing rings per layer are at
$z=0,\pm250,\pm550,\pm850,\pm1150\,\unit{\milli\metre}$, each
\qty{8}{\milli\metre} axially by \qty{3}{\milli\metre} radially.
A \qty{2}{\milli\metre}-wide beveled web at $u=\qty{14.5}{\milli\metre}$
joins a bare back-skin land to the ring tangent plane.

\begin{figure}[tbp]
\centering\includegraphics[width=\linewidth]{short-strip-stave}
\caption{Short-strip cold-stack and bearing schedule. Top: explanatory local
projection with direct thermal inserts, graphite spreader and insulated LV
rails; thin films and window cuts are specified in the module table. Bottom:
nine bearing stations and two independently served half-staves.}
\label{fig:short-stave}
\end{figure}

Each signed half-stave has a separate titanium U-loop, with legs at
$u=\pm\qty{12}{\milli\metre}$ and $w=-\qty{4.05}{\milli\metre}$.
The tube outside diameter is \qty{2.5}{\milli\metre}, wall thickness
\qty{0.14}{\milli\metre} and central return radius \qty{12}{\milli\metre}.
Straight runs span $|z|=16$--$1210\,\unit{\milli\metre}$.
Two backside copper rails have \qty{12}{\milli\metre} width per polarity and
\qty{0.17}{\milli\metre} thickness on \qty{0.15}{\milli\metre} polyimide.
HV, control and data retain separate patterned routing.

The $52\times30\times8\,\unit{\milli\metre^3}$ end-board allocations lie at
$z=\pm\qty{1230}{\milli\metre}$. Each half-stave uses one power/HV/control
bundle and two optical groups for its 14 modules. Each end therefore presents
288 power bundles, 576 optical groups and 288 cooling loops to twelve sector
collectors spanning $r=244$--$781.5\,\unit{\milli\metre}$ and
$|z|=1245$--$1310\,\unit{\milli\metre}$.

\subsection{Endcap discs and two-loop petals}
Six discs per end use the support datums in \cref{tab:strip-discs}. Twelve
\qty{30}{\degree} petals partition each disc. Five rings at
$289.5,381.25,473,564.75,656.5\,\unit{\milli\metre}$ carry
$48,60,84,96,108$ sensors, respectively: 33 per petal, 396 per disc and
\num{4752} in both endcaps. The sampled full-azimuth target annulus is
$r=243.5$--$703\,\unit{\milli\metre}$. The finite active outer corner reaches
\qty{704.909}{\milli\metre}, while the body corner reaches
\qty{708.424}{\milli\metre}. The physical active gap to the fixed long-strip
endcap is \qty{82.335}{\milli\metre}. These finite dimensions are distinct
from the sampled target; neither establishes reconstruction efficiency.

\begin{table}[!htbp]
\centering\small
\caption{Positive strip-disc support datums. Negative-end placements use the
corresponding signed datums. Values are rounded for display.}
\label{tab:strip-discs}
\begin{tabular}{@{}rrr@{}}\toprule
Disc & Short strips $|z|$ [mm] & Long strips $|z|$ [mm]\\\midrule
\csname @@input\endcsname generated/strip-disc-rows.tex
\bottomrule\end{tabular}
\end{table}

Alternating ring and azimuth elevations are $0,1.5,3,4.5\,\unit{\milli\metre}$.
The cooled plate spans $r=236$--$710\,\unit{\milli\metre}$ with a
\qty{0.6}{\milli\metre} inner seam gap. Its sandwich thicknesses match the
barrel. Inner and outer CFRP frames occupy $r=236$--$250$ and
$689$--$715\,\unit{\milli\metre}$; the outer frame retains
\qty{26}{\milli\metre} radial width. Peripheral board mounts transfer load
outside the outer sensor bodies. The $52\times20\times8\,\unit{\milli\metre^3}$
board is centred at $r=\qty{695}{\milli\metre}$ and local
$w=\qty{14}{\milli\metre}$, above the front fanout.

\begin{figure}[tbp]
\centering\includegraphics[width=\linewidth]{short-strip-endcap}
\caption{Short-strip endcap plan and one petal, with module elevations projected.
The inner and outer cooling circuits serve 16 and 17 modules. Blue and red
curves distinguish their centrelines; the two backside service buses are shown
schematically. Pipe contacts and local thermal windows are explicit in the
simulation model.}
\label{fig:short-endcap}
\end{figure}

Two independent loops per petal serve the inner three and outer two rings.
Their circular legs span $3$--$27\,\unit{\degree}$; returns use
\qty{10}{\milli\metre}-radius quarter-bends joined by radial straights.
Four separate rim ports lie at $r=\qty{712}{\milli\metre}$.
Tube diameters and walls match the barrel. Three protected 11-module power
islands use \qty{4}{\milli\metre} copper width per polarity each, retaining
\qty{24}{\milli\metre} total width and \qty{0.17}{\milli\metre} thickness.
The front fan spans $r=690$--$781.5\,\unit{\milli\metre}$ and
$w=20$--$100\,\unit{\milli\metre}$. Per end, the discs add 216 power bundles,
216 optical groups and 144 loops. Combined short-strip totals per end are
504 power bundles, 792 optical groups and 432 loops.

\subsection{Power, cooling and material boundaries}
Each finished power/HV/control bundle is allocated \qty{8}{\milli\metre}
outer diameter with two \qty{10}{\milli\metre\squared} LV conductors, independent
module HV pairs, control and shield. The protected distribution is
\qty{12}{\volt}, with at most 14 modules per power feed and a
\qty{11}{\volt} delivered-voltage bound. Optical groups retain the
12-module limit and \qty{3.6}{\milli\metre} package diameter. The shared trunk
occupies $r=717$--$781.5\,\unit{\milli\metre}$ for
$|z|=1310$--$3500\,\unit{\milli\metre}$, with \qty{0.5}{\milli\metre}
tray walls. The package dimensions require a manufacturable harness, grounding,
insulation, protection and connector design; no traffic reduction follows
merely from the cell-count reduction.

The nominal/adverse electronic power budgets are 13.88/19.83 W per module,
with 2 W sensor leakage and 5 W per physical end board. Coolant is required to
remain at or below \qty{-35}{\celsius}. Nominal barrel/endcap flow is 4/5 g/s
per leaf and adverse flow is 7 g/s. Bulk feed/return pairs retain 8/12 mm
outside and 6/10 mm inside diameters. The complete conditional adverse thermal
screen, including ASIC bonding, upstream sensor path and distributed leakage,
has only \qty{0.333}{\kelvin} margin to the \qty{-20}{\celsius} sensor goal.
Cold material/contact measurements, realistic power, pressure drop, dry-out,
flow balancing and structural qualification remain necessary.

Effective remote-service cells conserve the specified conductor and insulating
volumes plus residual void. Explicit local constituents include sensor, guard,
interconnect, readout fixtures, adhesives, graphite, TIM, backing, support,
pipes and coolant. The changed-component mass proxy gives approximately
\qty{114.6}{\kilogram} reduction when taking no cable-insulation saving and no
ASIC mass-reduction credit. This scoped design ledger includes added silicon
and support/contact allowances; it is not a measured tracker mass or a complete
directional radiation-length map. Current material and truth-fit results are
given in \cref{sec:tracker-evidence}; full reconstruction and engineering
qualification remain separate.
